\chapter{预备知识}

低空物流任务中的多无人机系统同时涉及飞行器动力学、姿态与轨迹控制以及基于环境交互的策略学习。这类问题的分析与求解涉及状态变量、控制输入、姿态表示和序贯决策模型等统一表述。本章包括四旋翼建模基础、基于 SE(3) 的四旋翼几何控制基础以及强化学习与连续马尔可夫决策过程基础。相关内容构成后续路径代价建模、轨迹跟踪控制和协同运输策略学习的理论基础，并统一全文中的变量体系与控制量表述。

\section{四旋翼无人机建模基础}

四旋翼无人机由机体、四个旋翼和动力系统构成，其运动由旋翼产生的推力与反扭矩共同驱动。相较于固定翼飞行器，四旋翼具有结构紧凑、悬停能力强和低速机动性能好的特点，因而在低空巡检、航拍测绘、环境监测和物资投送等任务中得到了广泛应用。四旋翼的控制与规划问题以刚体运动模型为基础。状态变量定义、坐标系设置、姿态表示方式以及控制输入与机体运动之间的关系如下。

四旋翼采用四个旋翼在机体平面内对称分布的构型，相邻旋翼旋向相反，对置旋翼旋向相同。该布局在产生总推力的同时抵消主要的旋翼反扭矩，悬停状态对应零合力矩平衡。四个旋翼同步增大转速时，总推力增大并形成上升运动；同步减小转速时，总推力减小并形成下降运动；前后两组旋翼转速不平衡时产生俯仰力矩；左右两组旋翼转速不平衡时产生滚转力矩；两组反向旋转旋翼之间存在转速差时产生偏航力矩。四旋翼的水平运动对应姿态倾斜后的推力水平分量，平动与转动始终耦合。

四旋翼在三维空间中的运动状态采用惯性坐标系和机体坐标系描述。惯性坐标系记为 $\Sigma_W=\{O_W,\mathbf e_1,\mathbf e_2,\mathbf e_3\}$，其中 $\mathbf e_1$ 和 $\mathbf e_2$ 为水平面内的两个正交单位基向量，$\mathbf e_3=[0,0,1]^\top$ 为竖直向上的单位基向量。该坐标系固定于地面，位置、速度、参考轨迹和重力方向均定义在该坐标系下。机体坐标系记为 $\Sigma_B=\{O_B,\mathbf b_1,\mathbf b_2,\mathbf b_3\}$，其原点位于四旋翼质心处，$\mathbf b_1$ 与 $\mathbf b_2$ 位于旋翼平面内，$\mathbf b_3$ 与旋翼平面法向一致并给出总推力方向。该坐标系随机体姿态变化，机体姿态、角速度和控制力矩均定义在该坐标系下。四旋翼的位置、速度、姿态和角速度分别记为 $p\in\mathbb R^3$、$v\in\mathbb R^3$、$R$ 和 $\Omega\in\mathbb R^3$，其中姿态矩阵 $R$ 属于特殊正交群（Special Orthogonal Group，SO(3)），系统状态写为 $x=(p,v,R,\Omega)$。

姿态表示采用欧拉角和旋转矩阵两种形式。欧拉角由滚转角、俯仰角和偏航角三个参数构成，分别记为 $\phi$、$\theta$ 和 $\psi$。按常用的 $ZYX$ 旋转顺序，机体姿态写为
\begin{equation}
    R=R_z(\psi)R_y(\theta)R_x(\phi),
\end{equation}
其中，偏航角 $\psi$ 表示机体绕竖直方向的航向变化，俯仰角 $\theta$ 表示机体绕侧向轴的抬头或低头运动，滚转角 $\phi$ 表示机体绕前向轴的侧倾运动。欧拉角直接对应机体的航向与倾斜状态；当俯仰角接近 $\pm \frac{\pi}{2}$ 时，参数化出现奇异性。后续姿态变量统一采用旋转矩阵 $R$。旋转矩阵 $R$ 描述机体坐标系相对于惯性坐标系的方向关系，并满足 $R^\top R=I$ 与 $\det(R)=1$。该表示在整个姿态空间内保持统一形式，适用于后续几何控制器的推导。对于任意向量 $x=[x_1,x_2,x_3]^\top\in\mathbb R^3$，记其对应的反对称矩阵为
\begin{equation}
    \hat{x}=
    \begin{bmatrix}
        0 & -x_3 & x_2\\
        x_3 & 0 & -x_1\\
        -x_2 & x_1 & 0
    \end{bmatrix},
\end{equation}
则有 $\hat{x}y=x\times y$，该映射记号用于姿态运动学表达。

四旋翼运动分为质心平动和绕质心转动两部分。其平动运动学和姿态运动学分别为
\begin{equation}
    \dot p=v,\qquad \dot R=R\hat{\Omega}.
\end{equation}
控制输入经推力和力矩进入平动加速度与角加速度方程，动力学关系如下。

设四旋翼质量为 $m$，惯性矩阵为 $J\in\mathbb R^{3\times 3}$，总推力大小为 $f$，机体系合力矩为 $\tau\in\mathbb R^3$，外部扰动分别记为 $d_f$ 和 $d_\tau$。四旋翼的平动动力学和转动动力学分别写为
\begin{equation}
    m\dot v=fR\mathbf e_3-mg\mathbf e_3+d_f,
\end{equation}
\begin{equation}
    J\dot\Omega+\Omega\times J\Omega=\tau+d_\tau.
\end{equation}
式中，$\Omega\times J\Omega$ 为刚体转动中的陀螺耦合项。平动方程描述质心受力，转动方程描述刚体转动，二者共同构成四旋翼的六自由度刚体动力学模型。

四旋翼的控制输入由四个电机驱动的旋翼产生。记第 $i$ 个电机的角速度为 $\omega_i$，旋翼旋向符号为 $\sigma_i\in\{-1,1\}$，单个旋翼产生的升力和反扭矩近似写为
\begin{equation}
    f_i=c_f\omega_i^2,\qquad m_i=\sigma_i c_m\omega_i^2,
\end{equation}
其中 $c_f$ 和 $c_m$ 分别为升力系数和反扭矩系数。记电机转速平方向量为 $s=[\omega_1^2,\omega_2^2,\omega_3^2,\omega_4^2]^\top$，旋翼推力向量为 $u=[f_1,f_2,f_3,f_4]^\top$，则广义控制输入与电机角速度平方、旋翼推力之间满足
\begin{equation}
    [\,f,\tau_x,\tau_y,\tau_z\,]^\top=A_\omega s=A[\,f_1,f_2,f_3,f_4\,]^\top,
\end{equation}
其中 $A_\omega$ 和 $A$ 由旋翼布局、机臂长度和旋转方向决定。所用四旋翼模型保留到电机角速度平方与单桨升力、反扭矩之间的静态映射层面。上述惯性坐标系与机体坐标系的相对关系如图~\ref{fig:coordinate_system} 所示。

\begin{figure}[htbp]
    \centering
    \includegraphics[width=0.6\linewidth]{pic/part0/orientation.jpg}
    \caption{坐标系示意图}
    \label{fig:coordinate_system}
\end{figure}

四旋翼属于典型的欠驱动系统。其在三维空间中具有位置和姿态共六个自由度，但独立控制输入只有总推力和三个力矩共四个量。四旋翼无法同时对三个平动方向和三个转动方向直接施加独立控制，而是通过姿态调整改变推力方向，再由推力投影实现横向与纵向机动。这一结构决定了平动控制与姿态控制相互耦合，轨迹跟踪和协同控制均建立在统一的刚体运动模型之上。

\section{基于 SE(3) 的四旋翼几何控制基础}

后续章节采用基于 SE(3) 的几何控制器作为四旋翼底层控制器。该控制器接收连续参考轨迹给出的位置信息、速度信息和加速度信息，也接收高层控制模块输出的期望合力与偏航参考。第四章的连续轨迹跟踪和第五章多智能体强化学习策略生成的高层控制信号均在该框架下转化为执行控制量。误差变量、期望姿态和控制律定义如下。四旋翼的运动状态由位置和姿态组成，其中位置属于欧氏空间 $\mathbb R^3$，姿态属于旋转群 SO(3)，整体状态定义在特殊欧氏群 SE(3) 上。几何控制直接在流形结构上构造误差和控制律，避免了欧拉角奇异性与姿态表示切换带来的局部不连续问题。

记四旋翼的期望轨迹为 $(p_d(t),v_d(t),R_d(t),\Omega_d(t))$。实际状态与期望状态之间的位置误差和速度误差定义为
\begin{equation}
    e_p=p-p_d,\qquad
    e_v=v-v_d.
\end{equation}
这两个误差量表征四旋翼质心在惯性空间中的跟踪偏差。

姿态误差和角速度误差分别定义为
\begin{equation}
    e_R=\frac{1}{2}\left(R_d^\top R-R^\top R_d\right)^\vee,
\end{equation}
\begin{equation}
    e_\Omega=\Omega-R^\top R_d\Omega_d.
\end{equation}
其中 $(\cdot)^\vee$ 表示 $\hat{(\cdot)}$ 映射的逆运算。误差 $e_R$ 表示当前姿态相对于期望姿态的几何偏差，$e_\Omega$ 表示当前角速度相对于期望角速度的机体系差异。该误差定义与旋转群结构保持一致。

平动控制参考向量定义为
\begin{equation}
    f_c=-k_p e_p-k_v e_v+mg\mathbf e_3+m\dot v_d,
\end{equation}
其中 $k_p>0$ 和 $k_v>0$ 为位置与速度反馈增益。向量 $f_c$ 包括位置误差反馈、速度误差反馈、重力补偿和期望加速度前馈。机体总推力沿 $\mathbf b_3$ 方向施加，期望姿态的第三轴与期望合力方向一致。

记期望偏航角为 $\psi_d$，水平参考方向为 $h_d=[\cos\psi_d,\sin\psi_d,0]^\top$，期望姿态矩阵按下式构造
\begin{equation}
    b_3^d=\frac{f_c}{\|f_c\|},
\end{equation}
\begin{equation}
    b_2^d=\frac{b_3^d\times h_d}{\|b_3^d\times h_d\|},
\end{equation}
\begin{equation}
    b_1^d=b_2^d\times b_3^d,
\end{equation}
\begin{equation}
    R_d=[\,b_1^d,\ b_2^d,\ b_3^d\,].
\end{equation}
$b_3^d$ 与期望合力方向一致，$h_d$ 约束偏航自由度，$R_d$ 为完整的期望姿态矩阵。

当控制输入为连续轨迹参考时，参考轨迹给出 $p_d(t)$、$v_d(t)$、$\dot v_d(t)$ 以及偏航参考 $\psi_d(t)$ 或等价姿态参考，控制器首先由 $e_p$、$e_v$ 和 $\dot v_d$ 计算平动控制参考向量 $f_c$，再结合 $\psi_d$ 构造 $b_3^d$、$b_2^d$、$b_1^d$ 和 $R_d$，随后由姿态误差与角速度误差计算总推力和控制力矩。当控制输入为高层力学参考时，高层模块直接给出期望合力向量及偏航参考，控制器不再通过位置误差与速度误差生成 $f_c$，而是将该期望合力向量直接作为姿态构造输入，按照上述方式确定 $b_3^d$ 和 $R_d$，再利用同一姿态误差定义和控制律求得推力与力矩指令。两类参考在输入来源上不同，但在姿态构造、误差计算和控制分配环节采用相同的控制器结构。

在得到误差量和期望姿态之后，几何控制器的总推力和控制力矩分别写为
\begin{equation}
    f=f_c^\top R\mathbf e_3,
\end{equation}
\begin{equation}
    \tau=-k_R e_R-k_\Omega e_\Omega+\Omega\times J\Omega
    -J\left(\hat\Omega R^\top R_d\Omega_d-R^\top R_d\dot\Omega_d\right).
\end{equation}
其中 $k_R>0$ 和 $k_\Omega>0$ 分别为姿态反馈增益和角速度反馈增益。总推力对应平动参考向量在当前姿态下的投影，控制力矩驱动姿态误差和角速度误差收敛。

四旋翼的最终执行器为四个电机，总推力和机体系力矩经控制分配映射为各旋翼推力或转速命令。记旋翼推力向量为 $u=[f_1,f_2,f_3,f_4]^\top$，则有
\begin{equation}
    u=A^{-1}[\,f,\tau_x,\tau_y,\tau_z\,]^\top.
\end{equation}
实际控制分配中还包含推力饱和、转速上界以及执行器约束投影等处理。整个控制链路覆盖轨迹参考、高层控制信号、姿态力矩与推力指令以及电机层执行量。

\section{连续马尔可夫决策过程基础}

强化学习研究智能体在环境交互条件下的序贯决策问题。其监督信息来自交互过程中获得的奖励信号，而非预先给定的输入输出标签。策略参数依据长期累积回报持续更新。在控制与决策问题中，强化学习框架由状态感知、动作选择、环境反馈和策略改进组成。

强化学习的基本要素包括智能体、环境、状态、动作、奖励与策略。智能体执行决策，环境响应动作并产生状态转移。状态 $s_t$ 描述时刻 $t$ 的环境信息，动作 $a_t$ 表示当前状态下的控制决策，奖励 $r_t$ 表示该时刻动作对任务目标的即时贡献，策略 $\pi$ 表示状态到动作的映射关系。在随机策略情形下，策略写为条件概率分布 $\pi(a|s)$；在确定性策略情形下，策略写为状态到动作的函数映射。强化学习的目标是通过优化策略提高整个任务过程的累积回报，而不是仅优化单步奖励。

若环境满足马尔可夫性质，即系统下一时刻的状态分布仅由当前状态和当前动作决定，而与更早时刻的历史信息无关，则强化学习问题表示为马尔可夫决策过程。一个马尔可夫决策过程（Markov Decision Process，MDP）记为 $\mathcal M=\langle \mathcal S,\mathcal A,\mathcal P,r,\gamma\rangle$，其中 $\mathcal S$ 为状态空间，$\mathcal A$ 为动作空间，$\mathcal P(s_{t+1}|s_t,a_t)$ 为状态转移概率，$r(s_t,a_t)$ 为奖励函数，$\gamma\in(0,1]$ 为折扣因子。对于给定策略 $\pi$，智能体与环境在每一时刻按照
\begin{equation}
    a_t\sim\pi(\cdot|s_t),\qquad
    s_{t+1}\sim\mathcal P(\cdot|s_t,a_t)
\end{equation}
的方式完成交互并获得奖励。状态转移序列和奖励序列组成强化学习中的样本轨迹，智能体、动作、环境反馈与状态更新之间的关系如图~\ref{fig:mdp_process} 所示。

\begin{figure}[htbp]
    \centering
    \includegraphics[width=0.8\linewidth]{pic/part0/RLbasic.png}
    \caption{马尔科夫决策过程示意图}
    \label{fig:mdp_process}
\end{figure}

当状态变量和动作变量均取连续值时，强化学习问题表示为连续 MDP。设 $s_t\in\mathcal S\subseteq\mathbb R^{n_s}$、$a_t\in\mathcal A\subseteq\mathbb R^{n_a}$，则系统在连续状态空间和连续动作空间中演化。无人机控制问题中的位置、速度、姿态和角速度属于连续状态变量，推力、姿态参考和力矩属于连续动作变量，连续 MDP 与飞行控制和协同决策问题相匹配。将从初始时刻到终止时刻的交互序列记为 $\tau=(s_0,a_0,r_0,s_1,a_1,r_1,\dots,s_T)$，则折扣累积回报写为 $G_0=\sum_{t=0}^{T-1}\gamma^t r_t$。强化学习的基本优化目标是寻找参数化策略 $\pi_\theta$，使其期望回报 $J(\pi_\theta)=\mathbb E_{\tau\sim\pi_\theta}[G_0]$ 尽可能大。

实际训练过程中，强化学习采用基于采样的方式近似求解上述优化问题。给定当前策略参数 $\theta$ 后，智能体在环境中执行若干轮交互，获得样本数据集
\begin{equation}
    \mathcal D=
    \left\{
    (s_t,a_t,r_t,s_{t+1})
    \right\}_{t=0}^{T-1}.
\end{equation}
样本数据集记录了策略在当前环境中的真实行为结果，也是策略参数更新的直接依据。复杂环境下的状态转移和奖励分布往往难以解析表达，基于样本的学习方式据此替代精确环境模型求解，并以实际交互数据推进策略性能改进。

强化学习训练由策略初始化、环境采样、样本评估和策略更新构成循环迭代。训练起始于初始策略，智能体依据当前策略在环境中执行动作并收集状态转移与奖励序列。采样结束后，利用所得样本估计当前策略在任务中的表现，并据此更新策略参数。随着采样与更新持续进行，策略由初始随机行为逐步收敛为具有明确任务导向的控制规律。对于连续控制问题，这种“交互采样---性能评估---策略改进”的训练闭环构成强化学习求解决策问题的基本框架。

\section{本章小结}

本章给出四旋翼控制与强化学习相关的预备知识。围绕四旋翼系统，统一定义了位置、速度、姿态和角速度等状态变量，建立了四旋翼的运动学与动力学模型，并写明了控制输入与执行器之间的映射关系。在此基础上，进一步定义了基于 SE(3) 的四旋翼几何控制基本形式，包括平动误差、姿态误差以及总推力和控制力矩的构造方式。随后概述了强化学习的基本要素、连续 MDP 形式以及基于采样的基本训练过程。上述内容构成后续模型建立、控制器设计与决策过程描述的理论基础。
